\begin{pastexam}[label={px:ncku-111-cce-ds-06}]{成大電通所 111 學年度 第 6 題}
（簡答）

(20 \textbf{pts}) 簡答題:
\begin{enumerate}[label=\alph*.]
\item (5 \textbf{pts}) For any nonempty binary tree, $T$, if $n_0$ is the number of leaf nodes and
$n_2$ is the number of nodes of degree 2. Pleases show the relation between $n_0$ and $n_2$ and
prove it.
\item (5 \textbf{pts}) Given the following Keys 62, 104, 67, 22, 103 and a hash table with 10
buckets, where each bucket has one slot. Please write down the result using the hash function
F(x) = (x mod 10). It adopts open addressing when a collision happens according to the quadratic
residue as follows:
\[
i_n=\begin{cases}
\left(F(x)+\left(\frac{n+1}{2}\right)^2\right) mod\ 10 & ,if\ n \text{ is odd}\\[4pt]
\left(F(x)-\left(\frac{n}{2}\right)^2\right) mod\ 10 & ,if\ n \text{ is even}
\end{cases}
\]
where $i_0=F(x)$ and $n$ denotes the collision number.
\begin{center}
\begin{tabular}{|c|c|c|c|c|c|c|c|c|c|}
\multicolumn{1}{c}{0} & \multicolumn{1}{c}{1} & \multicolumn{1}{c}{2} & \multicolumn{1}{c}{3} &
\multicolumn{1}{c}{4} & \multicolumn{1}{c}{5} & \multicolumn{1}{c}{6} & \multicolumn{1}{c}{7} &
\multicolumn{1}{c}{8} & \multicolumn{1}{c}{9}\\ \hline
\phantom{00} & \phantom{00} & \phantom{00} & \phantom{00} & \phantom{00} & \phantom{00} &
\phantom{00} & \phantom{00} & \phantom{00} & \phantom{00}\\ \hline
\end{tabular}
\end{center}
\end{enumerate}

The following figure describes a flow assignment in a flow network G(V, E). The $a/b$ notation in
each edge respectively denotes the current flow $a$ and capacity $b$ of the edge.
\begin{enumerate}[label=\alph*.]
\setcounter{enumi}{2}
\item (5 \textbf{pts}) Please show the residual network of the graph. Which algorithm can be used
to find an augmenting path in the residual network and what is its time complexity?
\item (5 \textbf{pts}) Please find the minimum cut of a network of the graph and what is the
maximum flow from S to T?
\end{enumerate}
\begin{center}
\begin{tikzpicture}[v/.style={circle,draw,inner sep=1pt,minimum size=6mm},l/.style={inner sep=1.5pt,font=\small},>={Stealth[length=2mm]}]
\node[v] (S) at (0,0) {S}; \node[v] (x) at (1.8,1.2) {x}; \node[v] (y) at (1.8,-1.2) {y};
\node[v] (z) at (4.4,1.2) {z}; \node[v] (P) at (4.4,-1.2) {P}; \node[v] (T) at (6.2,0) {T};
\draw[->] (S) -- node[l,above left] {2/4} (x);
\draw[->] (S) -- node[l,below left] {6/6} (y);
\draw[->] (x) -- node[l,above] {5/5} (z);
\draw[->] (x) -- node[l,left] {3/7} (y);
\draw[->] (y) -- node[l,pos=0.42,above left] {0/2} (z);
\draw[->] (z) -- (P);
\draw[->] (z) -- node[l,above right] {4/6} (T);
\draw[->] (y) -- node[l,below] {3/3} (P);
\draw[->] (P) -- node[l,below right] {4/4} (T);
\end{tikzpicture}
\end{center}
\end{pastexam}

\begin{pastsoln}
答案：(a) $n_0 = n_2 + 1$；(b) 第 2、3、4、7、9 格依序是 62、22、104、67、103；(c)(d) 見下（採用的
解讀下最大流 $10$，最小割容量 $10$）。 (a) 屬資結書第~7~章，(b) 屬第~12~章，(c)(d) 屬第~26~章，
本題放在第~26~章。

\textbf{(a)}（資結書第~7.1~節定理~7.2）\textbf{$n_0 = n_2 + 1$。}

\textbf{證明。} 設 $n_1$ 是度數 $1$ 的節點數，總節點數 $n = n_0 + n_1 + n_2$。 再數邊：除了根，
每個節點恰有一條連到父節點的邊，所以邊數是 $n - 1$；從父節點這邊數，度數 $i$ 的節點貢獻 $i$
條邊，邊數也等於 $n_1 + 2n_2$。 所以 $n_0 + n_1 + n_2 - 1 = n_1 + 2n_2$，化簡得
$n_0 = n_2 + 1$。 \hfill$\blacksquare$

\textbf{(b) 二次探測。} 第 $n$ 次碰撞後的位置依序是 $F,\ F + 1,\ F - 1,\ F + 4,\ F - 4,\ \ldots$
（模 $10$）；$F - k^2$ 是負數時取 $\bmod$ 的非負值（例如 $-1 \bmod 10 = 9$）。
\begin{center}
\begin{tabular}{cccl}
\toprule
鍵值 & $F(x)$ & 探測的位置 & 結果 \\
\midrule
62 & 2 & 2 & 放第 2 格 \\
104 & 4 & 4 & 放第 4 格 \\
67 & 7 & 7 & 放第 7 格 \\
22 & 2 & 2（滿），$2 + 1 = 3$ & 放第 3 格 \\
103 & 3 & 3, 4, 2, 7 都滿，$3 - 4 = -1 \equiv 9$ & 放第 9 格 \\
\bottomrule
\end{tabular}
\end{center}
\begin{center}
\begin{tabular}{|c|c|c|c|c|c|c|c|c|c|}
\multicolumn{1}{c}{0} & \multicolumn{1}{c}{1} & \multicolumn{1}{c}{2} & \multicolumn{1}{c}{3} &
\multicolumn{1}{c}{4} & \multicolumn{1}{c}{5} & \multicolumn{1}{c}{6} & \multicolumn{1}{c}{7} &
\multicolumn{1}{c}{8} & \multicolumn{1}{c}{9}\\ \hline
\phantom{00} & \phantom{00} & 62 & 22 & 104 & \phantom{00} & \phantom{00} & 67 & \phantom{00} & 103\\ \hline
\end{tabular}
\end{center}
$103$ 依序試了 $i_0 = 3$、$i_1 = 3 + 1 = 4$、$i_2 = 3 - 1 = 2$、$i_3 = 3 + 4 = 7$，都已經有資料，碰撞
$4$ 次後 $i_4 = (3 - 4) \bmod 10 = 9$ 才是空的。 這種 $\pm k^2$ 交錯的探測見資結書第~12.3~節。

\textbf{(c)(d) 原圖的問題。} 照圖上的方向（x$\to$y），給的流不滿足流量守恆：x 流入 $2$、流出
$5 + 3 = 8$；y 流入 $6 + 3 = 9$、流出 $0 + 3 = 3$。 另外 z$\to$P 這條邊沒有標數字。 如果 x、y 之間
的邊其實是 \textbf{y$\to$x}（流量 $3$、容量 $7$），而 z$\to$P 的流量是 $1$，每個點都守恆：
x 流入 $2 + 3 = 5$ = 流出 $5$；y 流入 $6$ = 流出 $3 + 0 + 3$；z 流入 $5 + 0$ = 流出 $4 + 1$；P 流入
$1 + 3$ = 流出 $4$；流量值是 $8$。 題目明說這是一個 flow assignment，所以下面採用這個解讀：
x--y 的邊是 y$\to$x，z$\to$P 的流量 $1$、容量記為 $c$（$c \ge 1$，原圖沒給，不影響答案）。 最後也
附上照原圖方向的答案。

\textbf{(c) 剩餘網路}（資結書第~26.2~節）。 每條邊 $(u, v)$ 流量 $f$、容量 $b$：剩餘容量
$b - f > 0$ 時畫順向邊 $u \to v$，$f > 0$ 時畫反向邊 $v \to u$（容量 $f$）。
\begin{center}
\begin{tikzpicture}[v/.style={circle,draw,inner sep=1pt,minimum size=6mm},l/.style={inner sep=1.5pt,font=\small},>={Stealth[length=2mm]},x=1.25cm,y=1.1cm]
\node[v] (S) at (0,0) {S}; \node[v] (x) at (1.8,1.2) {x}; \node[v] (y) at (1.8,-1.2) {y};
\node[v] (z) at (4.4,1.2) {z}; \node[v] (P) at (4.4,-1.2) {P}; \node[v] (T) at (6.2,0) {T};
\draw[->] (S) to[bend left=18] node[l,above left] {2} (x);
\draw[->] (x) to[bend left=18] node[l,below right] {2} (S);
\draw[->] (y) -- node[l,below left] {6} (S);
\draw[->] (z) -- node[l,above] {5} (x);
\draw[->] (y) to[bend left=18] node[l,left] {4} (x);
\draw[->] (x) to[bend left=18] node[l,right] {3} (y);
\draw[->] (y) -- node[l,pos=0.42,above left] {2} (z);
\draw[->] (z) to[bend left=18] node[l,pos=0.6,right] {$c-1$} (P);
\draw[->] (P) to[bend left=18] node[l,left] {1} (z);
\draw[->] (z) to[bend left=18] node[l,above right] {2} (T);
\draw[->] (T) to[bend left=18] node[l,pos=0.22,below=1pt] {4} (z);
\draw[->] (P) -- node[l,below] {3} (y);
\draw[->] (T) -- node[l,below right] {4} (P);
\end{tikzpicture}
\end{center}
（$c = 1$ 時沒有 z$\to$P 那條邊。）

找增廣路徑就是在剩餘網路上找一條 S 到 T 的路徑，用 \textbf{BFS}（或 DFS），時間 $O(V + E)$，
這張圖 $V = 6$。 每次都用 BFS 找最短的增廣路徑就是 \textbf{Edmonds--Karp} 演算法，最多增廣
$O(VE)$ 次，總時間 $O(VE^2)$（資結書第~26.4~節定理~26.3）。 這張剩餘網路上 BFS 找到
S$\to$x$\to$y$\to$z$\to$T，瓶頸 $\min(2, 3, 2, 2) = 2$；其中 x$\to$y 是反向邊，意思是把原本
y$\to$x 的 $3$ 單位退回 $2$ 單位。 增廣後流量值 $8 + 2 = 10$。

\textbf{(d) 最小割與最大流。} 增廣後 S$\to$x（$4/4$）、S$\to$y（$6/6$）都滿了，剩餘網路裡 S 走不
出去，可達集合是 $\{S\}$，所以最小割是
\[
(\{S\},\ \{x, y, z, P, T\}),\qquad \text{容量} = c(S, x) + c(S, y) = 4 + 6 = 10.
\]
由最大流最小割定理（資結書定理~26.6），最大流是 $\mathbf{10}$。 最小割不唯一：
$(\{S, x, y\}, \{z, P, T\})$ 的容量 $5 + 2 + 3 = 10$，$(\{S, x, y, z, P\}, \{T\})$ 的容量
$6 + 4 = 10$，也都是最小割。 增廣後的流量值 $10$ 等於這些割的容量，流不能再大、割也不能再小。

\textbf{照原圖方向（x$\to$y）的答案。} 這時 y 只能從 y$\to$z（$2$）、y$\to$P（$3$）送出，x 送給 y
的流也只會擠在這兩條邊上。 割 $(\{S, y\}, \{x, z, P, T\})$ 切到 S$\to$x、y$\to$z、y$\to$P
（x$\to$y 是從 T 側指回 S 側，不算），容量 $4 + 2 + 3 = 9$；而 S$\to$x$\to$z$\to$T 送 $4$、
S$\to$y$\to$z$\to$T 送 $2$、S$\to$y$\to$P$\to$T 送 $3$ 是一個值 $9$ 的流，所以最大流是 $9$，這個
割是最小割。 作答時把採用的方向寫清楚。
\end{pastsoln}
